\subsection{对应的性质 I}

\ssszh 定义方面的讨论既已交代完毕，现在转入正题，看看 $\mathbb{L}$ 应当是什么样子。更精确地说，我们列出两侧各自的一些基本对象，它们预期会互相对应。再加上与抛物诱导的相容性，这预期能在很大程度上确定该等价。

讨论的结构如下。我们在图 \ref{f:table} 中列出了互相对应的范畴。对图中每一行，我们会花一点篇幅描述两侧出现的对象，并说明为什么应当期待它们对应。最后我们看看，在这些情形中匹配缠绕算子意味着什么；特别是，会遇到前面许诺过的对 \eqref{e:loclangnaive} 的修正。



\begin{table}[htbp]
\centering
\begin{tabular}{|c|c|}
\hline
$LG$ 的范畴表示 & $\LS$ 上的范畴层 \\
\hline
$\Dmod(\on{Gr_G})$ & $\QCoh( \on{pt}/\cG)$ \\
$\Dmod(\Fl^{\on{aff}}_G)$ & $\QCoh( \widetilde{{\sN}}/\cG)$ \\
$\Dmod(LG/LN, \psi)$ & $\QCoh(\LS)$ \\
$\gkc\mod$ & $\QCoh(\Op)$ \\
\hline
\end{tabular}
\caption{局部几何朗兰兹下（猜想地）互相交换的一些对象。}
\label{f:table}
\end{table}


\ssszh 对图 \ref{f:table} 中的前两个例子，回忆 $p$-adic 表示论的一些基本特征会有助于理解。由光滑性，$p$-adic 群的任意表示 $\pi$ 对所有充分小的紧开子群 $K$ 都含非平凡的不变量 $\pi^K$。对这样的 $K$，若 $\pi$ 不可约，则不动向量 $\pi^K$ 是 Hecke 代数 $\on{Dist}(K \backslash G(\Fqt) / K)$ 的不可约表示，后者通过卷积自然地作用于它。这样便得到 $G(\Fqt)$ 的含 $K$-不变向量的不可约表示与该 Hecke 代数的不可约表示之间的双射。

当子群 $K$ 缩小时，相应的 Hecke 代数变大，这给出表示复杂度的一个基本度量。即：对给定的不可约 $\pi$，可以问使它有不平凡不变量的最大 $K$ 是什么；从某种意义上说，$K$ 越大，$\pi$ 越不复杂。


在局部朗兰兹的另一侧，$\Fqt$ 的绝对 Galois 群（更准确地说，其 Weil--Deligne 变体）带有典范的滤过，由递降的正规子群给出，称作分歧滤过。对给定的 $\cG(k)$-值 Galois 表示，可以问它的作用经过该滤过的哪个商，以此度量其复杂度。一条基本而引人注目的组织原则是：这两个复杂度度量在局部朗兰兹下互换。\footnote{更精确地说，表示的深度（由 Moy--Prasad 定义的某个非负有理数 \cite{moyprasad}）应当等于相应 Galois 表示的斜率，后者由上编号的分歧滤过度量。}


\subsection{对应的性质 II：非分歧表示}

\ssszh 由上述讨论可知，$p$-adic 群最容易理解的表示，应当是那些关于极大紧开子群 $G(\Fq[[t]])$ 有不动向量的表示。\footnote{应当指出，对 $GL_n$ 以外的群，还存在其他不共轭的极大紧开子群。}类似地，最简单的 Galois 表示应当是那些作用经过分歧滤过最后一个商的表示，即可以延拓到非穿孔圆盘上的 $\check{G}$-局部系统。具体地说，这些正是经过剩余域的 Galois 群分解的表示
%
\begin{equation} \label{e:galunramified}
 1 \rightarrow I \rightarrow \on{Gal}( \Fqt^s / \Fqt) \rightarrow
 \on{Gal}(\overline{\mathbb{F}}_q /
 \Fq) \simeq \widehat{\mathbb{Z}} \rightarrow 1,
\end{equation}
%
其中上标「$s$」表示可分闭包。

由于实际使用的是 Weil 群（按定义它把 $\widehat{\mathbb{Z}}$ 换成 $\mathbb{Z}$），非分歧 Galois 表示在同构意义下由 $\check{G}$ 中共轭类中的元素给出。也就是说，我们追踪 $1 \in \mathbb{Z}$ 在映射 $\mathbb{Z} \rightarrow \cG$ 下的像。在这个最简单的例子中，局部朗兰兹归结为球 Hecke 代数与 $\cG$ 的特征标环之间的 Satake 同构，即
%
\[    \on{Dist}(G(\Fq[[t]]) \backslash G(\Fqt) / G(\Fq[[t]])) \simeq
\on{Fun}(\cG(k) / \cG(k)),
\]
%
其中右端表示 $\cG$ 上伴随不变的多项式函数，把 $\cG$ 看作 $\on{Spec} k$ 上的簇 \cite{satake63}。过渡到特征标，就得到：$G(\Fqt)$ 的不可约非分歧表示由 $\cG(k)$ 中的（半单）共轭类参数化，正如所期望的那样。

\begin{re} 对初次接触这些想法的读者，我们要强调：甚至上面出现的 Hecke 代数的交换性都不是先验显然的。这恰好说明，过渡到关于一个大紧子群的双不变分布，会多么戏剧性地简化问题。\end{re}

\ssszh 为几何版本作准备，我们应当指出：每个不可约非分歧表示都唯一地（相差一个标量）嵌入万有非分歧表示
%
\begin{equation} \label{e:fungr}
	    \Fun( G(\Fqt) /G(\Fq[[t]])).
\end{equation}
%
而且，要从万有情形提取出每一个，只需取 \eqref{e:fungr} 中在球 Hecke 代数的右卷积下按相应特征标变换的子模。



最后，我们讨论上述内容在局部几何朗兰兹中的类比。在自守侧，把仿射 Grassmann 簇上的 D-模范畴
%
\[
	   \Dmod(\Gr_G) = \Dmod(LG/L^+G)
\]
%
当作上述万有非分歧表示的替代物；读者现在应当已经熟悉这种类比了。在谱侧，穿孔圆盘上任何能延拓到整个圆盘的 de Rham 局部系统都是可平凡化的，于是得到叠点
%
\[
	    \on{pt}/\cG \rightarrow \LS
\]
%
相应地把 $\QCoh( \on{pt}/\cG)$ 看作 $\LS$ 上的范畴层。


还应当指出：不仅有万有情形的类比，也有单个不可约非分歧表示的类比。更详细的讨论留到第 \ref{ss:locaff} 节。

\begin{re}应当指出，非分歧情形虽是最基本的，但它在算术朗兰兹与几何朗兰兹中都完全根本。即在整体朗兰兹对应中（这里不详述），adelic 群 $G(\mathbb{A}_F)$ 的任意不可约光滑表示都分解为各位置处 $p$-adic 群 $G(F_x)$ 的不可约表示的限制张量积，而除有限多个位置外都是非分歧的。为了哪怕理解这样的无限张量积，都必须实质性地用到非分歧表示有典范的生成直线，即其球不变量。在整体几何理论中，迄今已知的几乎一切都只是关于处处非分歧（或许还有野分歧）自守表示的类比。
\end{re}
